\section{A resource-preserving static routing sparsifier}
\label{sec:routing-helper}

The recursive cycle construction needs a sparse set of representative edges
through which all edges of an auxiliary graph can be routed.  Its application
also associates resources with the original edges: using a representative
edge consumes its resources in either direction, and repeated traversals must
be charged repeatedly.  We prove a static sparsification statement that
preserves these loads.  The support bound depends on the number of incident
labels, even when the number of demands is much larger.

Throughout this section, let $S=2^s\ge2$ and $M=2^u$ satisfy
$S^{1/a}\le M\le S^2$, where $a\ge10$ is an integer and $u>0$ is divisible
by $2a^2$.  Define
\begin{equation}\label{route:parameters}
 b=M^{1/a^2},\qquad
 \Lambda=2^{1000a^2}(1+u),\qquad P_*=\Lambda^{400a^2},
 \qquad P_*\le b.
\end{equation}
The last inequality is a checked prerequisite.  In particular, $b$ is an
integer.  Powers of $\Lambda$ are logarithmic powers of the master size for
each fixed $a$.

\begin{lemma}[Static routing with resources]\label{lem:routing-helper}
Let $F$ be an explicitly listed undirected multigraph with identified edge
copies, and let $v_+$ be its number of incident vertex labels.  For each
edge $f$, an explicit list specifies resource counts
$m_{xf}\in\{1,\ldots,10\}$; unlisted counts are zero.  Suppose that
\[
 \deg_F(v)\le C_0,\qquad \sum_{f\in F}m_{xf}\le C_0
 \quad\text{for every vertex $v$ and resource $x$},
\]
where $C_0\ge1$ is an integer, degrees count incidences, and both $2|F|$
and the total resource-list length are at most $M^2$.
Assume that identifiers and numerical input data have
$(\log S)^{O_a(1)}$ bits.  There is a uniform randomized algorithm that
either reports detected failure or returns a set $F'\subseteq F$ and,
for every $f=vw\in F$, an explicit walk $Q_f$ from $v$ to $w$ in $F'$,
such that
\begin{equation}\label{route:guarantees}
 |F'|\le P_*v_+,\qquad
 |Q_f|\le P_*,\qquad
 \sum_{f\in F}\ \sum_{e\text{ traversed in }Q_f}m_{xe}
       \le P_*C_0\quad\text{for every $x$}.
\end{equation}
The inner sum counts traversal multiplicities and ignores orientation.
A loop demand may be assigned the empty walk.  Every returned result is
verified.  For each fixed reliability parameter $\sigma\ge0$, the time
on every computation path is
\begin{equation}\label{route:time}
 (1+\text{input list size})\,b\,M^{o_a(1)},
\end{equation}
and the probability of detected failure is at most
$2^{-\sigma}S^{-800}$.  The same statement holds conditionally when the
input is determined by earlier computation and this call uses fresh bits.
\end{lemma}

We build the representative set in stages, each routing at least half
the remaining demands.  An attempt samples edges according to their
endpoint degrees and gives each retained edge its inverse-probability
weight.  Introduce a distinct \emph{stub} for each edge incidence.
A flow between two equal halves of the stubs either identifies a cut
for recursion or pairs almost all the stubs by walks in the sampled
graph.  The sampling estimate ensures that the cuts are sparse with
high probability.  We defer their crossing demands to later stages
and recurse on the two induced pieces with fresh samples.  When almost
all stubs are paired, completing the unpaired stubs with temporary fake
edges gives a perfect matching on the stubs.  A piece that completes the
prescribed rounds without a further cut obtains an overlay made from these matchings;
the median rule makes it expanding with high probability.  Each nonfake
overlay edge stores its sampled-edge walk.

We then route the demand stub pairs through this overlay.  Partition its
labels into groups and send each demand packet on a short random walk
conditioned to reach its target's group.  Additional independent walks,
starting and ending within each group, supply the edges of smaller
graphs on which routing continues.  Recursion reaches singleton groups,
and substituting the saved walks gives overlay routes with bounded
length and congestion.  Discard routes that touch a fake edge or an
excessively long real embedding.  Congestion bounds the number of lost
demands, while the sampling weights bound resource use after the
remaining real embeddings are substituted.  The final verification
checks these walks and their loads before accepting the stage.

We use the following parameters:
\begin{equation}\label{route:internal-parameters}
\begin{gathered}
 r_m=\Lambda^3,\qquad K_s=\Lambda^{230a^2},\qquad
 \chi=\Lambda^{110a^2},\qquad \epsilon_0=\Lambda^{-100a^2},\\
 D_p=\Lambda^{105a^2},\qquad H_*=\Lambda^{70a^2},\qquad
 L_0=\Lambda^{200a^2},\qquad t=\Lambda^{10},\qquad q_v=\Lambda^3.
\end{gathered}
\end{equation}
Here $r_m$ counts overlay rounds, $K_s$ sets the sampling scale, $\chi$
multiplies flow capacities, $\epsilon_0$ sets the deficiency scale, and
$D_p$ caps the cut depth.  The bounds $H_*$ and $L_0$ control overlay routes
and retained real embeddings; $t$ is half the conditioned-walk length.
The symbol $q_v$ prescribes the same number of virtual jobs at each vertex,
a fixed logarithmic power for fixed $a$; it is unrelated to the accuracy
parameter $q$ of the priced engine.

For an unweighted multigraph $H$, let $\partial_H X$ be the edge-copy cut
of $X$, and let $\vol_H(X)=\sum_{v\in X}\deg_H(v)$.
Parallel copies count separately; loops count twice in volume and never
in a cut.  We suppress the subscript when the graph is clear.  Its graph
conductance, whenever there is a cut with two positive volumes, is
\[
 \phi(H)=\min_{\substack{X\subseteq V(H)\\
                 \vol_H(X)>0,\ \vol_H(V(H)\setminus X)>0}}
 \frac{|\partial_H X|}
      {\min\{\vol_H(X),\vol_H(V(H)\setminus X)\}}.
\]
We invoke conductance and mixing only on nonsingleton active walk graphs;
singleton router instances are settled directly.  For the lazy walk used
below, the corresponding transition-kernel conductance is $\phi(H)/2$.

\subsection{Stages, pieces, and normalized sampling}

Remove loop demands by giving them empty walks.  Work on the remaining
demands in stages.  One successful stage must route at least half its
input demands while using, separately for support, maximum hop count,
and resource load, the bound $\Lambda^{350a^2}$ in place of $P_*$ in
\eqref{route:guarantees}.  All these properties, including endpoints and
selected edge identities, are checked before a stage is accepted.  A
stage gets at most
\begin{equation}\label{route:attempts}
 T=2(\sigma+1000\log_2 S+10)
\end{equation}
fresh attempts (rounding up if necessary).  Exhausting these attempts is
a detected failure.  There are at most $2u+2$ successful stages, since
$|F|\le M^2/2$ and the residual number halves each time.

An attempt processes a recursive family of pieces.  A piece is an edge
list induced by the vertex side of each earlier cut; cut edges are
discarded for this stage.  Empty lists need no further processing.
Pieces at a fixed depth have disjoint edge lists and disjoint incident
vertex sets.  Process pieces by depth, using fresh samples at each
depth.  We cap the recursion depth at $D_p$ and abort an attempt
that would exceed the cap.  Let a nonempty piece have $s'$ edges and
positive original degrees $d_i$.  Retain each edge $e=ij$ independently
with probability $p_e$, the smallest number of the form $2^{-j}$,
$j\ge0$, that is at least
\[
 \min\{1,K_s(1/d_i+1/d_j)\}.
\]
A retained edge receives integral weight $\omega_e=1/p_e$.
Before any flow call, require
\begin{equation}\label{route:sample-checks}
 \sum_{e\text{ retained}}\omega_e\le2s',\qquad
 \#\{e\text{ retained}\}\le4K_s\#\{i:d_i>0\}.
\end{equation}
A violation aborts this attempt.  Volumes in the following sampling
estimate use the original piece degrees, not the sampled degrees.

\begin{lemma}[Simultaneous cuts and resource loads]
\label{route:sampling}
For a sampled piece, except with probability at most
$4s'\exp[-K_s/(100\chi^2)]$, every vertex set $X$ satisfies
\begin{equation}\label{route:cut-approximation}
 \left|\sum_{e\in\partial X\text{ retained}}\omega_e
                -|\partial X|\right|
 \le\frac{\min\{\vol(X),\vol(\bar X)\}}{4\chi}.
\end{equation}
This event implies the total-weight check in
\eqref{route:sample-checks}; the sample-count check also holds with
exponentially high probability in $K_s$.  Conditional on the input
pieces at a fixed depth, for each resource $x$,
\begin{equation}\label{route:sample-resources}
 \sum_{\substack{e\text{ retained}\\\text{over all pieces at this depth}}}
          \omega_e m_{xe}\le2C_0
\end{equation}
except with probability at most $\exp(-K_s/30)$.
\end{lemma}

\begin{proof}
Put $D=\operatorname{diag}(d_i)$ on the nonisolates and set
\[
 A_e=D^{-1/2}(\mathbf e_i-\mathbf e_j)
             (\mathbf e_i-\mathbf e_j)^\mathsf{T}D^{-1/2}.
\]
These are positive semidefinite rank-one matrices,
$\|A_e\|=1/d_i+1/d_j$, and $\|\sum_e A_e\|\le2$.
If $\xi_e$ is the retention indicator, write
$Y_e=(\xi_e/p_e-1)A_e$.  Deterministically retained edges contribute
zero.  For any other edge, $p_e\ge K_s\|A_e\|$, whence
\[
 \|Y_e\|\le K_s^{-1},\qquad
 \E Y_e^2=(p_e^{-1}-1)A_e^2\preceq K_s^{-1}A_e.
\]
Thus the variance parameter is at most $2/K_s$.  The self-adjoint
matrix Bernstein inequality, applied to both signs
\cite[Theorem~6.1(ii)]{TroppBernstein}, states that independent centered
self-adjoint summands of norm at most $J$ and total variance norm at
most $V$ satisfy
\[
 \Pr\left\{\left\|\sum_eY_e\right\|\ge z\right\}
 \le2d\exp\left(-\frac{z^2}{2(V+Jz/3)}\right),
\]
where $d$ is the matrix dimension.  Here $d\le2s'$, $J=1/K_s$,
$V\le2/K_s$, and $z=1/(4\chi)$ give the claimed probability bound.
Testing the quadratic form on $D^{1/2}\mathbf1_X$ gives additive
error $\vol(X)/(4\chi)$ in its cut weight.  Testing the complement
gives the other bound in \eqref{route:cut-approximation}.  Testing a
singleton bounds its weighted degree by
$(1+1/(4\chi))d_i$; summing degrees proves the weight check.

Dyadic rounding gives
\[
 \sum_ep_e\le2K_s\sum_{ij=e}(1/d_i+1/d_j)
       =2K_s\#\{i:d_i>0\}.
\]
The scalar Chernoff inequality therefore proves the sample-count
check with failure at most $\exp(-K_s\#\{i:d_i>0\}/3)$.
If the threshold exceeds the total possible edge count, failure is
of course impossible.

For \eqref{route:sample-resources}, condition on all lists at this
depth.  They are disjoint sublists of the input, so the expected sum
for $x$ is at most $C_0$.  A nondeterministic summand
$\xi_e m_{xe}/p_e$ lies in $[0,10C_0/K_s]$, since piece degrees are at
most $C_0$ and $1/p_e\le C_0/K_s$.  After centering, its absolute
value has the same upper bound.  The sum of variances is at most
$10C_0^2/K_s$, by bounding each second moment by its maximum value
times its expectation.  Deterministic summands have zero variance.
For an upward deviation of $C_0$, scalar Bernstein consequently
gives
\[
 \Pr\{\text{sum}>2C_0\}
 \le\exp\left(
 -\frac{C_0^2}{2(10C_0^2/K_s+10C_0^2/(3K_s))}\right)
 \le e^{-K_s/30}.
\]
\end{proof}

\subsection{Median cuts, integral matchings, and an expanding overlay}

The median-projection and matching-overlay argument has the cut-matching
structure of Khandekar, Rao, and
Vazirani~\cite[Sections~3--3.2]{KhandekarRaoVazirani2006}.
Here the vertices are incidence stubs, and the proof separately controls
deficient flow, resource loads, exact arithmetic and rejected attempts.

The piece has $N_p=2s'$ distinct stubs, one for each original edge
incidence.  Run at most $r_m$ rounds.  Each round divides the stubs
into two equal halves by the median rule below.  Let $l_i$ and $h_i$
count the left and right stubs at original vertex $i$.  Form a flow
network with source capacities $l_i$, sink capacities $h_i$, and
both directed arcs of capacity $\chi\omega_e$ for every retained
edge.  Its maximum possible value is $s'$.

If its deficiency $\Delta=s'-\text{maximum value}$ is at least
$\epsilon_0s'$, use the source-side vertex set $X$ of a minimum cut
to split the piece and stop its rounds.  The cut formula gives
\begin{equation}\label{route:deficiency-cut}
 l(X)-h(X)=\Delta+
             \chi\sum_{e\in\partial X\text{ retained}}\omega_e.
\end{equation}
The left side equals $h(\bar X)-l(\bar X)$, so both original volumes
are at least $\epsilon_0s'$.  It is also at most their minimum.
On the event in Lemma~\ref{route:sampling},
\begin{equation}\label{route:piece-cut}
 |\partial X|\le
 \frac{2\min\{\vol(X),\vol(\bar X)\}}{\chi}.
\end{equation}
Indeed the weighted cut is at most the minimum volume divided by
$\chi$, and \eqref{route:cut-approximation} adds only one quarter
of that upper bound.  Each child's internal edge count is at most
$(1-\epsilon_0/2)s'$, because the opposite side has volume at least
$\epsilon_0s'$.  A nonempty branch thus has fewer than
$1+4\epsilon_0^{-1}\log M<D_p$ cuts.  Moreover,
$\min\{\vol(X),\vol(\bar X)\}\le s'$, so summing
\eqref{route:piece-cut} at each depth shows that all cuts discard
at most the fraction
\begin{equation}\label{route:cut-loss}
 2(D_p+1)/\chi
\end{equation}
of the attempt's initial demands.  This charging uses disjoint
piece lists at each depth; it does not charge each demand to every
piece at that depth.

When $\Delta<\epsilon_0s'$, extract the integral flow as individually
identified pairs of left and right stubs with explicit sampled-edge
walks.  Pair all remaining left and right stubs arbitrarily and call
these pairs \emph{fake edges}.  The result is a perfect matching on
the stubs.  Save this matching as one round of an overlay multigraph,
together with real embeddings for its nonfake edges.

These flow calls use Lemma~\ref{lem:integral-flow}.  Padding to
$O((\chi+1)s')$ makes every capacity polynomially bounded in the
padded size, since the weight check precedes the call.  Alternatively,
Lemma~\ref{lem:capacity-scaling} provides the same fixed primitive
when local capacities are presented with longer bit strings.  Path
extraction has a separate explicit charge: pair incoming and outgoing
flow units at each internal vertex, also pairing the units on source
and sink arcs with their stubs, and trace from source units.  Remaining
cycles are omitted.  There are at most
\begin{equation}\label{route:flow-throughput}
 2\chi\sum_{e\text{ retained}}\omega_e\le4\chi s'
\end{equation}
interior arc units per round.  Lists of these units and their pairings
therefore construct all paths within this bound plus the network size.
No uncharged expansion of a large integral flow is used.

\begin{lemma}[The median rule]\label{route:overlay}
There is a median rule for which a piece completing all $r_m$ rounds
has an $r_m$-regular overlay of conductance at least $1/(4r_m)$,
except with probability at most
$4N_p^3\exp[-r_m/(80\Lambda)]$.  This guarantee holds for any choice
of perfect matching across each chosen pair of halves.
\end{lemma}

\begin{proof}
For a perfect matching let $A$ be its averaging matrix: on each
matched pair it replaces two coordinates by their average.  Start
with $P_0=I$, and after round $j$ set
$P_j=A_jP_{j-1}$.  All these matrices are doubly stochastic.  At a
round's start, let $p_i$ be the current rows and choose a fresh vector
$\zeta$ of independent uniform signs.  Order the real numbers
$\langle p_i,\zeta\rangle$, resolving ties by stub identifiers, and
put their first and last halves on opposite sides.

Write $\bar p=\mathbf1/N_p$ and
$\Psi=\sum_i\|p_i-\bar p\|_2^2$.  The centered projected values
$z_i=\langle p_i-\bar p,\zeta\rangle$ have mean zero over $i$, and
$Z=\sum_i z_i^2$ satisfies
\[
 \E Z=\Psi,\qquad \E Z^2\le3\Psi^2.
\]
For the second inequality, sign fourth moments give
$\E\langle v,\zeta\rangle^4\le3\|v\|_2^4$, and
Cauchy--Schwarz bounds each cross term in $Z^2$.
The elementary second-moment inequality now gives
$\Pr\{Z\ge\Psi/2\}\ge1/12$ when $\Psi>0$.
The sign tail estimate and a union bound over pairs give, with
failure probability at most $2N_p^2e^{-\Lambda/2}$,
\[
 |\langle p_i-p_j,\zeta\rangle|^2
       \le\Lambda\|p_i-p_j\|_2^2\quad\text{for every $i,j$}.
\]

Choose a median value $m$ between the two central centered values $z_i$.
For any pair $i,j$ on opposite sides,
$(z_i-z_j)^2\ge(z_i-m)^2+(z_j-m)^2$.  Every stub occurs once in
the matching, and the sum of squared distances is minimized at the
mean, zero.  Consequently, for \emph{every} matching across the
halves,
\[
 \sum_{ij\text{ matched}}(z_i-z_j)^2
       \ge\sum_i(z_i-m)^2\ge Z.
\]
The exact decrease under averaging is
$\frac12\sum_{ij\text{ matched}}\|p_i-p_j\|_2^2$.
With conditional probability at least $1/20$ this is at least
$\Psi/(4\Lambda)$.  The parameter values ensure that the sign-tail
failure above is smaller than $1/12-1/20$.  If $\Psi=0$, it stays
zero.  Thus
\[
 \E\Psi_{r_m}\le N_p\exp[-r_m/(80\Lambda)].
\]
For this analysis one can continue with arbitrary median-crossing
matchings after a piece has been cut; hence stopping a piece early
does not introduce conditioning on successful mixing.  If any entry
of $P_{r_m}$ is below $1/(2N_p)$, then
$\Psi_{r_m}\ge1/(4N_p^2)$.  Markov's inequality proves the stated
failure bound.

It remains to identify which walk this product describes.  Since
$P_{r_m}=A_{r_m}\cdots A_1$, its row starting at stub $i$ describes
averaging steps in the order $r_m,r_m-1,\ldots,1$.  Place one unit
of mass at each stub of a set $X$ with $|X|\le N_p/2$, and run this
\emph{reverse-round} walk.  All product entries are at least
$1/(2N_p)$, so at least $|X|/4$ mass finishes outside $X$.
At every intermediate time, double stochasticity keeps mass at each
stub at most one.  A matching edge $\{v,w\}$ at its single time step
therefore carries total mass at most
$(\mathrm{mass}(v)+\mathrm{mass}(w))/2\le1$.  The number of overlay
edge occurrences crossing $X$ is at least $|X|/4$.  Every stub has
degree $r_m$, which proves conductance at least $1/(4r_m)$.

The construction does not form dense products.  To compute a new
projection, apply the earlier matching averages in chronological
order to $\zeta$, obtaining $P_{j-1}\zeta$.  This costs
$O(jN_p)$ elementary averages.  All entries are dyadic rationals
with denominator dividing $2^{j-1}$, so exact projections, comparisons,
and tie resolution use $O(r_m+\log N_p)$ bits per number.
\end{proof}

\subsection{Recursive routing on the overlay}

The next task is to route the demand stub pairs through the overlay
with bounded path length and total traversal load on each overlay
edge.  These two bounds will control the lengths and resource loads
after the saved real-edge embeddings are substituted.

The recursive sampling of virtual graphs uses the group-conditioned
walk construction of Ghaffari, Kuhn, and Su
\cite[Section~3.1.2]{GhaffariKuhnSu2017}.  We give the packet routing,
congestion, and finite-computation estimates needed here, allowing
nonuniform stationary distributions and counting repeated traversals.

Pair the original stubs according to their demand edges and orient
each pair arbitrarily, giving one packet with a specified source and
final target.  Both the initial sources and the final targets are
distinct.  At recursive calls current packet locations can coincide,
but their final targets remain distinct.  We route these packets
through the overlay, temporarily permitting fake edges.

Every walk graph retains distinct parallel edge copies and uses
incidence degrees: a loop has two incidences.  Its lazy walk holds
with probability $1/2$ and otherwise chooses a uniform incident
edge occurrence.  Holding incurs no traversal.  An actual traversal
of a loop is counted once.  The stationary law is
$\pi(v)=d_v/\sum_i d_i$.

We use the standard conductance-to-mixing argument for reversible
chains~\cite[Section~3]{SinclairJerrum1989}. The proof fixes our graph
conductance normalization and allows loops and parallel copies.

\begin{lemma}[Mixing from conductance]\label{route:cheeger}
In a connected undirected multigraph with at least two vertices and
conductance $\phi$, the
lazy walk has spectral gap at least $\phi^2/4$.  If the graph has at
most $N_p$ labels, degree ratio at most $D_c\le\Lambda^{1/20}$,
and conductance at least $1/(4r_m)$, then for every $h\ge t$,
\begin{equation}\label{route:pointwise-mixing}
 \tfrac12\pi(y)\le P^h(x,y)\le2\pi(y)
 \quad\text{for all $x,y$}.
\end{equation}
\end{lemma}

\begin{proof}
For any nonnegative vector $f$ supported on at most half the degree
volume, integrate the boundary sizes of the sets
$\{i:f_i^2>\xi\}$ over $\xi\ge0$.  The result is
\[
 \sum_{ij\text{ edge}}|f_i^2-f_j^2|
       \ge\phi\sum_i d_i f_i^2.
\]
Cauchy--Schwarz, together with
$\sum_{ij\text{ edge}}(f_i+f_j)^2\le2\sum_i d_i f_i^2$, implies
\begin{equation}\label{route:cheeger-energy}
 \sum_{ij\text{ edge}}(f_i-f_j)^2
       \ge\frac{\phi^2}{2}\sum_i d_i f_i^2.
\end{equation}
Both formulas remain valid with loops counted by their two incidences.
For an arbitrary vector $g$ of degree-weighted mean zero, subtract a
degree-weighted median $m$.  The positive and negative parts of
$g-m$ each have support volume at most half.  Their edge energies
sum to at most that of $g$, while their squared degree norms sum to
\[
 \sum_i d_i(g_i-m)^2
       =\sum_i d_i g_i^2+m^2\sum_i d_i
       \ge\sum_i d_i g_i^2.
\]
Applying \eqref{route:cheeger-energy} to both parts proves the
nonlazy gap bound $\phi^2/2$, by the Dirichlet quotient, and hence
the lazy bound $\phi^2/4$.  The nonlazy reversible operator is a
contraction in the degree-weighted Euclidean norm; making it lazy
places its spectrum in $[0,1]$.

Apply its contraction on mean-zero densities to a point mass.
Cauchy--Schwarz in the stationary inner product gives relative
pointwise error at most
\[
 \frac{\exp(-h\phi^2/4)}{\sqrt{\pi(x)\pi(y)}}
       \le\frac{\exp(-h\phi^2/4)}{\pi_{\min}}.
\]
Here $\pi_{\min}\ge1/(D_cN_p)$, $N_p\le M^2$,
$\phi\ge1/(4\Lambda^3)$, and $t=\Lambda^{10}$.
Thus the last expression is at most
$D_cM^2\exp(-\Lambda^4/64)<1/2$, establishing
\eqref{route:pointwise-mixing}.
\end{proof}

The grouped recursive virtual-job construction below adapts the
hierarchical routing idea of Ghaffari--Kuhn--Su~\cite[Lemma~3.1]{GhaffariKuhnSu2017}.
The sequential implementation, explicit caps, and resource-load bounds
needed here are established below; we do not invoke their distributed-algorithm
guarantee.

At a recursive call, a singleton label requires no routing.  Otherwise
divide its label set into $\min\{b,\#\text{labels}\}$ groups whose
sizes differ by at most one.  For each packet, draw a $2t$-step walk
from its present location conditioned to end in the group of its final
target.  For each vertex, also draw $q_v$ independent \emph{virtual
jobs} with the same length, starting there and conditioned to end in
that vertex's own group.  The start and endpoint of every virtual job
form one undirected edge within its group; its sampled walk is saved
as the edge's embedding.  Recurse on each of these virtual graphs
with the packets at their newly sampled locations and with their
original final targets.  On return, substitute the stored embeddings,
reversing them when necessary.

All job draws within a call are independent conditional on its input.
A conditioned draw is implemented by ordinary walk trials, retaining
the first whose endpoint is in the prescribed group.  Allow at most
$\Lambda^3b$ trials per job.  For a uniform choice among $d$ incident
occurrences, draw $\lceil\log_2 d\rceil$ fair bits and reject indices
outside $\{0,\ldots,d-1\}$; allow at most $\Lambda$ strings for
one such choice.  A cap violation aborts the attempt.  These rules
also cover graphs with loops and parallel edges.  The group-size rule
makes the maximum size after $j$ recursive levels at most
$\lceil N_p/b^j\rceil$.  Hence the recursion depth is at most
$2a^2$, since $b^{2a^2}=M^2\ge N_p$.

\begin{lemma}[Packet and virtual-job router]\label{route:router}
Conditional on an overlay satisfying Lemma~\ref{route:overlay}, the
procedure above produces explicit routes for all packets with both
maximum length and total traversal congestion on each overlay edge
at most $H_*$, except with probability less than $M^{-20}$ throughout
the checked parameter range.  All its data sizes and running
times remain bounded by the specified caps even when it fails.
\end{lemma}

\begin{proof}
First analyze ideal unlimited rejection, stopping the analysis when
an inductive good event fails.  The current degree ratio is allowed
to be $D_c$, initially one, with the allowance multiplied by $20$
per level.  Thus
\[
 D_c\le20^{2a^2}\le\Lambda^{1/20}.
\]
The inductive conductance lower bound is $1/(4r_m)$, so
Lemma~\ref{route:cheeger} applies.  A group $A$ occupies at least
$1/(2b)$ of the labels; hence $\pi(A)\ge1/(2D_cb)$.
An ordinary trial succeeds with probability at least
$1/(4D_cb)$.  The endpoint of a successful trial has distribution
\begin{equation}\label{route:conditioned-endpoints}
 \frac{P^{2t}(x,y)}{P^{2t}(x,A)}
 \in\left[\frac14,4\right]\frac{\pi(y)}{\pi(A)}
 \subseteq\left[\frac{1}{4D_c|A|},\frac{4D_c}{|A|}\right]
 \quad(y\in A).
\end{equation}
Distinct final targets mean that at most $|A|$ packets seek group
$A$.  Packet arrivals at a specified child vertex therefore have
expectation at most $4D_c$, and Chernoff gives an upper bound of
$q_v$ with overwhelmingly high probability.

Exactly $q_v|A|$ virtual jobs seek $A$.  If $I_z$ counts their
endpoints at $z$, then $\E I_z\le4D_cq_v$, and with overwhelmingly
high probability $I_z\le16D_cq_v$ for every $z$.  The child degree
is exactly $q_v+I_z$: a loop contributes its start incidence and
its endpoint incidence.  Thus every child degree is at least $q_v$
and at most $(1+16D_c)q_v\le20D_cq_v$, proving the claimed next
degree ratio.

For any $X\subseteq A$ with $1\le|X|\le|A|/2$, each of the
$q_v|X|$ jobs starting in $X$ leaves $X$ with probability at least
$1/(8D_c)$, by \eqref{route:conditioned-endpoints}.  Chernoff yields
\begin{equation}\label{route:child-cut}
 \#\{\text{virtual edges from $X$ to $A\setminus X$}\}
       \ge\frac{q_v|X|}{16D_c}
\end{equation}
except with probability $\exp[-q_v|X|/(64D_c)]$.
Union over subsets by cardinality costs at most
\begin{equation}\label{route:cut-union}
 \sum_{j\ge1}N_p^j\exp[-q_vj/(64D_c)].
\end{equation}
For an arbitrary nontrivial cut, apply \eqref{route:child-cut} to
its smaller-cardinality side $X$.  Its volume is at most
$20D_cq_v|X|$, and the smaller of the two cut volumes is no larger
than this volume.  The child conductance is at least
$1/(320D_c^2)\ge1/(4r_m)$.  This establishes all graph and packet
invariants required at the next level.

We next bound loads of \emph{all} jobs at one call before any
substitution.  Set $Q=q_v$, $W=\sum_i d_i$, and
$d_{\min}=\min_i d_i$.  At most $2Q$ jobs start at any vertex:
there are $Q$ virtual jobs and at most $Q$ packets.  At most
$(Q+1)|A|$ jobs seek group $A$, because virtual jobs seeking $A$
start in $A$ and packet final targets are distinct.

Fix an oriented incidence $x\to y$ of a specific edge.  At a step
in the first half, the remaining suffix has at least $t$ steps.
Conditioning a job on its prescribed terminal group increases the
probability of this step's traversal by at most four: the ratio is
$P^h(y,A)/P^{2t}(z,A)\le2\pi(A)/(\pi(A)/2)$, where $z$ is its
start.  Without conditioning, the aggregate initial density is at
most $2Q/\pi_{\min}=2QW/d_{\min}$ relative to stationarity, and
the walk preserves this bound.  A stationary lazy step traverses
this oriented incidence with probability $1/(2W)$.  Therefore its
expected conditioned load is at most
\begin{equation}\label{route:first-half-load}
 4\frac{2QW}{d_{\min}}\frac{1}{2W}
       =\frac{4Q}{d_{\min}}.
\end{equation}

In a second-half step, the prefix before traversal has at least
$t$ steps.  If $j$ steps remain after this incidence and the job
seeks $A$, its conditional traversal probability is at most
\[
 \frac{2\pi(x)\,(1/(2d_x))\,P^j(y,A)}{\pi(A)/2}
       =\frac{2P^j(y,A)}{W\pi(A)}.
\]
Summing first over jobs seeking $A$ and then over the groups gives
\begin{align}
 \sum_A\frac{2(Q+1)|A|}{W\pi(A)}P^j(y,A)
 &=2(Q+1)\sum_A\frac{|A|}{\vol(A)}P^j(y,A)\notag\\
 &\le\frac{2(Q+1)}{d_{\min}}\sum_AP^j(y,A)
 =\frac{2(Q+1)}{d_{\min}}.\label{route:second-half-load}
\end{align}
This calculation cancels the rare-conditioning factor group by
group.  In particular it introduces neither the number of groups
nor a degree-ratio factor.  Adding the two oriented incidences of
an undirected edge, including the two incidences of a loop, bounds
its expected traversal load at any step by $24q_v$.

Conditional on the current input, different job streams are
independent.  At one fixed step, their indicators of traversing
a fixed edge are independent Bernoulli variables, regardless of
dependence between different steps of one job.  Scalar Chernoff
therefore bounds this load by $\Lambda^6$ with overwhelmingly high
probability.  A union bound over edges and the $2t$ steps gives
total local edge load at most $2t\Lambda^6\le\Lambda^{20}$;
local walk lengths are also at most $\Lambda^{20}$.

To account for substitution, suppose child routes have maximum
length and edge congestion at most $Z$.  Initial packet walks
contribute at most $\Lambda^{20}$ load to a current edge.
Each virtual-edge embedding is then used at most $Z$ times, so
all substitutions contribute at most $Z\Lambda^{20}$.
The same inequality bounds length, using the maximum embedding
length.  Starting at the leaves and iterating through at most
$2a^2$ levels bounds both quantities by
\begin{equation}\label{route:substitution}
 (2\Lambda^{20})^{2a^2}\le H_*.
\end{equation}

For completeness, all probabilities just used are conditional on
the actual recursively specified input.  Endpoints chosen at a
parent create the child graph, but at that child the new streams
are fresh independent streams.  This permits induction up to the
first failed graph or packet invariant and avoids any assumption
that the child graph is independent of its packet starts.
The group-cut bound \eqref{route:cut-union}, all point-count and
edge-step tails, and all cap tails are below $M^{-40}$ with the
parameters above.  A conditioned job exceeds its trial cap with
probability at most
\[
 \left(1-\frac{1}{4D_cb}\right)^{\Lambda^3b}
       \le\exp[-\Lambda^3/(4D_c)],
\]
and a uniform incidence choice exceeds its string cap with
probability at most $2^{-\Lambda}$.  Coupling capped rejection to
unlimited rejection until the first cap justifies using these
bounds with the ideal distributions.

Finally, the bounds on the work of a failed run require no good
event.  Packets never duplicate; each vertex always initiates
exactly $q_v$ virtual jobs; group sizes shrink by the fixed rule;
and every trial and walk has an explicit cap.  Across a router
depth there are at most $q_vN_p$ virtual edges, with incidence
arrays for constant-index access.  Even if congestion bounds fail,
let $\mathcal L_h$ bound a fully expanded packet route with $h$ levels remaining.
The newly sampled prefix has at most $2t$ steps, and each edge of the
child route expands to at most $2t$ steps.  Thus
\[
 \mathcal L_0=0,\qquad \mathcal L_h\le2t(1+\mathcal L_{h-1})
       \le\sum_{j=1}^{h}(2t)^j\le2(2t)^h
       \qquad(1\le h\le2a^2).
\]
This is a fixed power of $\Lambda$.  These deterministic size bounds also bound
all elementary choice counts.  Under $P_*\le b$, a crude total
of pieces, calls, edges, steps, and capped choices in an entire
attempt is less than $M^{20}$.  A union bound thus proves the
asserted router probability, with ample margin.
\end{proof}

\subsection{Filtering real embeddings and verifying a stage}

The router's good event bounds route lengths and congestion in the
overlay, where fake edges are still allowed.  We must now discard
routes using fake edges or excessively long real embeddings.  The
congestion bound limits how many demands this discards and controls
the resource load of the remaining expanded routes.

At a completed piece, discard every packet route using a fake
overlay edge or an overlay edge whose real embedding has more than
$L_0$ sampled-edge hops.  This filtering is performed on overlay
identities before any long real embedding is substituted.  In one
piece there are fewer than $\epsilon_0s'r_m$ fake edges.  By
\eqref{route:flow-throughput}, the total length of all real
embeddings is at most $4r_m\chi s'$, so at most
$4r_m\chi s'/L_0$ overlay edges have a long real embedding.
Each such edge is used by at most $H_*$ packets, since traversal
congestion bounds the number of packets using it.  The fraction
discarded within completed pieces is therefore at most
\begin{equation}\label{route:filter-loss}
 H_*\epsilon_0r_m+\frac{4H_*r_m\chi}{L_0}.
\end{equation}
Adding \eqref{route:cut-loss}, the total lost fraction is at most
\[
 2(D_p+1)/\chi+
       H_*\epsilon_0r_m+4H_*r_m\chi/L_0<\tfrac12.
\]
Thus more than half the original stage demands survive all cuts
and filtering on the good events.

Replace each remaining overlay traversal by its real sampled-edge
embedding.  Its direction specifies whether to use the embedding
or its reverse.  The resulting walks have correct endpoints and
at most $H_*L_0$ hops.  In each round, an undirected sampled edge
$e$ appears at most $2\chi\omega_e$ times among the real embeddings,
by its two directed capacities.  It is therefore traversed at most
\begin{equation}\label{route:sampled-edge-load}
 2H_*r_m\chi\omega_e
\end{equation}
times by the final walks from its completed piece.

Completed pieces are vertex-disjoint even when they stop at
different depths: no completed piece has descendants, and sibling
pieces use disjoint vertex sets.  Their samples therefore contain
at most $4K_sv_+$ distinct original edges.  For resources, use
\eqref{route:sampled-edge-load} and apply
\eqref{route:sample-resources} separately at each of the at most
$D_p+1$ depths.  This bounds every resource load by
\begin{equation}\label{route:stage-load}
 4H_*r_m\chi(D_p+1)C_0.
\end{equation}
The parameters give each required stage margin explicitly:
\begin{equation}\label{route:stage-margins}
 4K_s\le\Lambda^{350a^2},\qquad
 H_*L_0=\Lambda^{270a^2}\le\Lambda^{350a^2},\qquad
 4H_*r_m\chi(D_p+1)\le\Lambda^{350a^2}.
\end{equation}

The algorithm accepts a stage only after checking these three
bounds with $\Lambda^{350a^2}$, checking that at least half the
demands have been assigned, and checking every returned walk's
edge identities and endpoint sequence.  It counts total
traversals of each selected edge and multiplies these counts into
the edge's sparse resource list to verify resource loads.  In
particular, neither expansion nor the probabilistic premises of
the construction need to be trusted by a subsequent user of the
sparsifier: accepted walks satisfy the deterministic interface.

\subsection{Failure probability and finite computation}

\begin{proof}[Proof of Lemma~\ref{lem:routing-helper}]
At a fixed cut depth, condition on its input lists before making
its fresh samples.  Lemma~\ref{route:sampling} applies to these
lists even though earlier cuts were random.  Likewise the median
potential and the recursive routing estimates are conditional on
their respective inputs.  The sum of all bad-event probabilities
in one attempt is less than $1/2$ in the checked parameter range.
Here are quantitative counts for this assertion.  There are at
most $M^2(D_p+1)$ pieces and at most $M^2$ resources having nonzero
input data.  The matrix and sample-count tails, resource tails,
and product-mixing tails displayed above are each below $M^{-40}$,
including their per-piece factors.  For the router, the group-cut
union is bounded by \eqref{route:cut-union}; the point-count and
edge-step tails use $q_v=\Lambda^3$ and threshold $\Lambda^6$;
the rejection tails are the ones proved in
Lemma~\ref{route:router}.  Each is below $M^{-40}$.  Their total
number, including capped random-choice positions, is less than
$M^{20}$.  These comparisons follow directly from
$\Lambda=2^{1000a^2}(1+u)$, $a\ge10$, and $P_*\le b$:
every polynomial count in the listed powers of $\Lambda$ is
absorbed by the latter inequality, while every displayed tail
has an exponent dominating a constant multiple of $u$.
All good events imply that the stage passes its verification.

After each unsuccessful attempt, the unchanged residual graph is
tried again with fresh bits.  Conditional failure probability is
at most $1/2$ per attempt; hence failure of all $T$ attempts has
probability at most $2^{-T}$.  Taking a union bound over at most
$2u+2$ stages, using $u\le2\log_2 S$, gives
\[
 (2u+2)2^{-T}\le2^{-\sigma}S^{-800}.
\]
This also proves the stated conditional guarantee for adaptive
inputs.  Taking the union of selected supports over accepted
stages and retaining each demand's first accepted walk yields
\eqref{route:guarantees}, because
\[
 (2u+2)\Lambda^{350a^2}\le P_*.
\]
Support and resource loads add over stages; maximum hop count
does not increase when supports are united.

We finish with the cost on every path, including rejected
attempts.  At any fixed cut depth, piece edge counts sum to at
most the starting edge count.  Sampling scans their lists once.
For each piece there are at most $r_m$ flow calls, each on padded
size $P=O((1+\chi)s')$.  Lemmas~\ref{lem:integral-flow}
and~\ref{lem:capacity-scaling} bound their exact finite-word cost
by padded size times $M^{o_a(1)}$, uniformly over these local
sizes.  In particular, $P\le O((1+\chi)M^2)$ and the flow
primitive's subpolynomial overhead at these sizes is
$M^{o_a(1)}$.  The total padding over a depth is linear in its
edge count times a fixed power of $\Lambda$.  The weight check
\eqref{route:sample-checks} bounds flow throughput by
\eqref{route:flow-throughput} even for a sample with inaccurate
cuts.  Constructing all integral paths is therefore charged
explicitly.  Projecting sign vectors requires only repeated
matching averages, $O(r_m^2N_p)$ operations per piece, and never
requires a dense matrix.  Sorting projected values to choose each
median costs $O(N_p\log N_p)$ comparisons and is absorbed in the
same logarithmic-power allowance.

At a completed piece, the initial overlay has $r_mN_p/2$ edges.
At each routing depth there are at most $q_vN_p$ virtual edges
and $O(q_vN_p)$ jobs across its calls.  Store all incident-edge
arrays explicitly, so a walk step does not scan a neighborhood.
A job uses at most $\Lambda^3b$ trials, each of $2t$ steps, and
each incident choice has at most $\Lambda$ bit-string trials.
These caps supply the sole additional nonlogarithmic factor $b$.
The deterministic recursion-depth and substitution bounds from
Lemma~\ref{route:router} control intermediate walk-list sizes
even on bad random outcomes.  Long real embeddings are filtered
before substitution; all remaining expansions have a further
factor of at most $L_0$.  Thus explicit output construction is
also linear in the input size times $b$ and a fixed power of
$\Lambda$.

Finally, sparse dictionaries translate original labels, edge
identifiers, and resource names to local indices and back.
Verification first counts traversals by selected edge and then
scans its resource list, so a large resource list is not recopied
at every traversal.  All comparisons, capacity operations,
resource-load products, and dyadic averages use bounded
multiword integers.  In particular, the projection denominator
has at most $r_m$ binary digits; local path lengths and load
counters have the deterministic size bounds just proved; and
all input bit strings have the assumed logarithmic-power width.
The normalized matrices, potentials, and stationary probabilities
are used only for analysis; the algorithm never computes their
square roots or eigenvectors.  The sampling probabilities are
found by exact integer comparisons, and an event of probability
$2^{-j}$ is drawn using $j$ fair bits.  Here $j\le2u$, since the
positive piece degrees are at most $M^2$.  Each random sign uses
one fair bit.  The ordinary multiword implementations of all
arithmetic operations add only logarithmic powers for fixed $a$.
Individual generated bits,
initialization of arrays, allocation, and movement of these
records are included in the same charges.  Multiplication by the
capped numbers of cut depths, routing depths, stages, and
attempts leaves precisely \eqref{route:time}.  Since all caps
are enforced independently of the good events, this is a bound
on every computation path.
\end{proof}
